% !TeX root = ../../../../../main.tex
% chapters/sections/chapter3/subsections/households/ces_general.tex

\subsection{消費指数と物価指数}
\label{sec:chapter3_households_indices}
本節では上記の2段階最適化問題で用いる各種の消費指数とそれに対応する物価指数を定義していく。
消費指数はCES（Constant Elasticity of Substitution）関数により定義され、
このCES関数は家計が多様な財の消費の組み合わせから得る効用をモデル化するものである。
つまりCES型の消費指数を引数とした効用関数を採用することにより
各家計の財が互いに幾分差別化されている独占的競争状態を表現することができる。
\begin{equation}
\label{form:chapter3_households_indices_ces_eq}
    U=\left(\alpha_1c_1^{\frac{\theta-1}{\theta}}+\alpha_2c_2^{\frac{\theta-1}{\theta}}\right)^{\frac{\theta}{\theta-1}}
\end{equation}
ここで\(\alpha\)は財への根源的な好みの度合い（バイアス）を示し\(\theta\)は
代替の弾力性を示すパラメータである。
内側の指数\(\frac{\theta-1}{\theta}\)の値は財の代替のしやすさを決定する。
指数が小さい（\(\theta\)が1に近い）場合、ある財の限界的な効用の減少が早いため、
各財に消費を分散させる方が効用が高まり代替が難しい状況となる。
逆に指数が大きい（\(\theta\)が大きい）場合、限界効用の減少が緩やかであるため各財は代替が容易な状況となる。
なお全体にかかっている外側の指数\(\frac{\theta}{\theta-1}\)は
内側の指数による数学的な変換を元に戻す調整項である。
これにより、集計量である消費指数が元の消費という経済的に意味のある単位を維持することが保証される。
具体的には、消費の単位を「単位」とすると
\begin{equation}
\label{form:chapter3_households_indices_ces_units}
    \left[(\text{単位})^{\frac{\theta-1}{\theta}}\right]^{\frac{\theta}{\theta-1}}=\text{単位}^{(\frac{\theta-1}{\theta})\cdot(\frac{\theta}{\theta-1})}=\text{単位}^1
\end{equation}
となり消費指数は元の消費と同じ次元をもつ。

このCES型関数から限界代替率（\(MRS_{1,2}\)）が次のように導かれる。
\begin{equation}
\label{form:chapter3_households_indices_mrs}
    MRS_{1,2}=\frac{\alpha_1}{\alpha_2}\left(\frac{c_1}{c_2}\right)^{-1/\theta}
\end{equation}

\(MRS_{1,2}\)は特定の消費時点において財1が財2の何倍の価値があるかを表すため、
1財の代替のしにくさを表す指標と考えることができる。
この\(MRS_{1,2}\)の水準は式の通り\(\alpha_1,\alpha_2,c_1,c_2,\theta\)によって決定される。
しかし\(\alpha_1,\alpha_2,\theta\)は家計の選好を表す固定されたパラメータである。
そこで\(\alpha\)の影響を取り除き、
変動する変数である消費の比率（\(c_2/c_1\)）と\(MRS_{1,2}\)の純粋な関係を分析する必要が生じる。

その際に変化率の関係、すなわち弾力性を考えることで2つの変数の関係を分析することができる。
具体的には\(MRS_{1,2}\)の変化の割合が消費の比率\(c_2/c_1\)の変化の割合の何倍かは一定値となり、
それが代替の弾力性\(\theta\)となる。

以下では、家計\(h\)の効用の源となる総消費指数を合成する。
この総消費指数の対数（log）をとったものが\(h\)が消費から得られる効用となる。
総消費指数を合成するため、まずそのもととなる自国財消費指数と外国財消費指数を定義する。
これらの消費指数はすべて先ほど説明したCES型関数の特殊なものである。
どのように特殊であるかといえば、
自国財および外国財消費指数は個別財の消費を\(\alpha\)を均一化（対称性を仮定）して
CES型関数により合成したものである。
また総消費指数は自国財および外国財消費指数を\(\theta=1\)としてCES型関数により合成したものである。
以上の設計思想は以下の表のようにまとめられる。

\begin{table}[H]
    \centering
    \caption{消費指数の設計思想の比較}
    \label{tab:chapter3_households_indices_philosophy}
    \begin{tabular}{lll}
        \toprule
        & 自国財および外国財消費指数(\(c_t^{h\to H},c_t^{h\to F}\)) & 総消費指数(\(c_t^{h\to W}\)) \\
        \midrule
        \(\alpha\)の扱い & 対称性を仮定するため不要 & 非対称性（ホームバイアス）のため残す \\
        \(\theta\)の扱い & パラメータとして残す & \(\theta=1\)に固定（コブ＝ダグラス型） \\
        分析の主役 & \(\theta\)(代替の弾力性) & \(\alpha^H\)(ホームバイアス) \\
        \bottomrule
    \end{tabular}
\end{table}

このモデルの設計が持つ意義は分析の焦点を明確に分離している点にある。
自国財および外国財消費指数においては個別財間の代替の弾力性\(\theta\)が
生産者の価格設定行動や物価の粘着性を左右する。
この\(\theta\)を分析の主役とするため各財好みを表す\(\alpha\)は平準化して消去する。
一方で総消費指数においては2国間の消費・貿易の基本的なパターン
（どれだけ輸入し、どれだけ国内で消費するか）を決定するのは
家計の根源的な内外財への選好の偏り、すなわちホームバイアス\(\alpha^H\)である。
代替の弾力性\(\theta\)は為替レートの変化などに対する反応の仕方を決めるが、
それは\(\alpha^H\)が設定した基本的な貿易パターン上での変動に過ぎない。
したがって国際的なパターンを分析する上では\(\alpha^H\)がより根源的なパラメータであると考え、
この\(\alpha\)を分析の主役とするため\(\theta=1\)という
「最も普通の」場合を考えて\(\theta\)を消去する。

以降の節でこれらの消費指数の厳密な定義とそこから導かれる関係式を詳述する。